\documentclass[ecta,nameyear,draft]{econsocart}
\usepackage{xr-hyper}
\usepackage{amsmath,amsfonts,amssymb,amsthm,graphicx,booktabs,array,longtable}
\usepackage{algorithm,algpseudocode,bm,enumitem}
\usepackage[T1]{fontenc}
\usepackage{mathpazo}
\usepackage{microtype}
\usepackage{needspace}
\RequirePackage[colorlinks,citecolor=blue,linkcolor=blue,urlcolor=blue]{hyperref}
\startlocaldefs
\newtheorem{theorem}{Theorem}
\newtheorem{proposition}[theorem]{Proposition}
\newtheorem{lemma}[theorem]{Lemma}
\newtheorem{corollary}[theorem]{Corollary}
\newtheorem{assumption}{Assumption}
\theoremstyle{definition}
\newtheorem{definition}{Definition}
\theoremstyle{remark}
\newtheorem{remark}{Remark}
\newtheorem{example}{Example}
\newcommand{\E}{\mathbb E}
\newcommand{\R}{\mathbb R}
\newcommand{\HH}{\mathcal H}
\newcommand{\LL}{\mathcal L}
\newcommand{\Act}{\mathcal A}
\newcommand{\sg}{\operatorname{sg}}
\newcommand{\tr}{\operatorname{tr}}
\newcommand{\Lip}{\operatorname{Lip}}
\DeclareMathOperator*{\argmax}{arg\,max}
\DeclareMathOperator*{\argmin}{arg\,min}
\endlocaldefs

\let\NBOLongtable\longtable
\def\longtable{\Needspace{7\baselineskip}\NBOLongtable}
\externaldocument{revisions/2026-10-05-r19-integrated/build/supp_refs}[supp.pdf]
\externaldocument{revisions/2026-10-05-r19-integrated/build/applications_refs}[applications.pdf]
\begin{document}
\begin{frontmatter}
\title{Neural Bellman Operators}
\runtitle{Neural Bellman Operators}
\begin{aug}
\author[id=au1,addressref={add1}]{\fnms{Qian}~\snm{QI}\ead[label=e1]{qiqian@pku.edu.cn}}
\address[id=add1]{Peking University}
\end{aug}
\begin{abstract}
This paper develops Neural Bellman Operators for policy evaluation and feasible
improvement in controlled economies. A common continuation connects current
choices to policy-specific future payoffs. We derive decision-loss bounds from
centered continuation errors, a transport bound for changes in future economic
primitives, and a true-objective certificate for continuous scalar decisions.
These results give explicit conditions for reuse and refresh while retaining
the boundary, utility-domain, and deviation requirements of the original
economic applications. The capital evidence includes protected high-dimensional
payoff comparisons, conditional prediction-risk assessment, and simultaneous
finite-catalogue bounds. A new prospective experiment varies one to 1,024
queries and compares seven complete procedures at a common economic tolerance,
charging fitting, querying, failed checks, and verification. NBO attains the
scalar-interval target in every declared stream; conventional surrogates remain
competitive and are not excluded from the cost frontier. The results distinguish
a reusable learned continuation from a claim of uniform neural dominance or
an unverified solution of the full diffusion control problem.
\end{abstract}
\begin{keyword}
\kwd{Dynamic programming}\kwd{Neural approximation}\kwd{Recursive utility}
\kwd{Endogenous preferences}\kwd{Viscosity solutions}\kwd{Dynamic games}
\end{keyword}
\end{frontmatter}
\input{revisions/2026-10-05-r19-integrated/manuscript/introduction.tex}
\input{revisions/2026-10-04-r16/manuscript/literature.tex}
\input{revisions/2026-10-05-r19-integrated/manuscript/literature_addition.tex}
\input{revisions/2026-10-05-r19-integrated/manuscript/targets.tex}
\input{revisions/2026-10-05-r19-integrated/manuscript/decision_main.tex}
\input{revisions/2026-10-04-r16/retained/manuscript/model.tex}
\input{revisions/2026-10-04-r16/retained/manuscript/method.tex}
\input{revisions/2026-10-04-r16/retained/manuscript/algorithm.tex}
\input{revisions/2026-10-04-r16/manuscript/theory_guide.tex}
\input{revisions/2026-10-04-r16/manuscript/menu_main_theory.tex}
\input{revisions/2026-10-04-r16/manuscript/new_theory_summary.tex}
\input{revisions/2026-10-04-r16/manuscript/computational_design.tex}
\input{revisions/2026-10-05-r18/manuscript/current_results.tex}
\input{revisions/2026-10-05-r18/manuscript/risk_main.tex}
\input{revisions/2026-10-05-r18/manuscript/catalogue_main.tex}
\input{revisions/2026-10-05-r19-integrated/manuscript/prospective_main.tex}
\input{revisions/2026-10-04-r16/manuscript/historical_evidence_summary.tex}
\input{revisions/2026-10-04-r16/retained/manuscript/economic_scope.tex}
\input{revisions/2026-10-05-r19-integrated/manuscript/conclusion.tex}
\bibliographystyle{ecta-fullname}
\bibliography{revision_reference,revisions/2026-10-05-r19-integrated/references}
\end{document}
